% =============================================================================
%  paper/supplementary.tex -- Digital Discovery supplementary material
%
%  Companion file to paper/main.tex.  Per Royal Society of Chemistry
%  conventions for Digital Discovery, the supplementary information is
%  hosted as a separate PDF on the journal's SI portal; this file is
%  \input{}-ed from main.tex after the appendices.
%
%  Conventions used throughout this file:
%    - Every numerical claim is flagged \MEASURED{} or \DESIGN{}; no
%      number is re-derived inside the supplementary body.  The SI
%      only restates, cross-references, and exposes the underlying
%      artefact.
%    - Internal reports under molmetal/reports/ are cited via the
%      footnote:* keys in paper/refs.bib.
%    - Local appendix files (closure_theorem.tex, betanf_semantics.tex)
%      are cross-referenced by \label, not re-printed here.
% =============================================================================

\section*{Supplementary material}
\addcontentsline{toc}{section}{Supplementary material}

\paragraph{Scope.}
This document collects the implementation, hyperparameter, per-pocket,
ablation, scoring-protocol, parity, normalisation and closure proof
details that underpin the manuscript but would interrupt its narrative.
Ten supplementary items (S1--S10) are included.  Each item states its
provenance explicitly and uses the same \MEASURED{}/\DESIGN{} markers as
the main text.

\paragraph{Honest framing.}
Every supplementary item below cites the corresponding internal report
or local appendix under \texttt{molmetal/reports/},
\texttt{molmetal/configs/}, \texttt{molmetal/molmetal\_lam/}, or
\texttt{paper/appendices/}.  No number in the main manuscript is
re-derived inside this supplementary material; SI sections only restate
and cross-reference the upstream artefact.  Two items (S3
\emph{Per-pocket results} and S10 \emph{Closure theorem + 5-click
algorithm}) are explicitly marked \DESIGN{} because the underlying
experiments / proofs are still in flight at the time of writing
(\texttt{footnote:round12\_pilot}, \texttt{footnote:round13\_sweep},
and \texttt{WF-Lambda-4}; see \texttt{TODO/pending/}).

% ---------------------------------------------------------------------------
% S1 -- Implementation details: 9-layer MLC architecture
% ---------------------------------------------------------------------------

\subsection*{S1. Implementation details: the 9-layer MLC architecture}
\label{supp:s1-mlc-impl}

\paragraph{Provenance.}
The Mol-Metal hybrid decomposes into nine layers of metrics, organised
into the 4-layer rewriting core (atoms, bonds, applications,
combinators) plus 5 observational layers (geometric priors, empirical
scoremix, lambda-only metrics, click-rule registry, closure-theorem
state).  Figure~1 of the manuscript (see also
\texttt{paper/figures/fig1\_mlc\_architecture.{png,svg,pdf}}) is the
canonical diagram and is reproduced from the WF-Paper-2 build
(\texttt{molmetal/reports/wf\_paper2\_figures.md}).

\paragraph{Layer catalogue.}
Layer numbers, metric names and the source module / report are listed
in \texttt{molmetal/reports/lambda\_layer\_metrics.md}.  The full
9-layer breakdown (substrate $\rightarrow$ typed-application $\rightarrow$
combinator $\rightarrow$ click-rule registry $\rightarrow$ metal-geometry
prior $\rightarrow$ empirical scoremix $\rightarrow$ lambda-only metrics
$\rightarrow$ MCTS-closure state $\rightarrow$ oracle-budget trace) is
exposed in that report.  Each layer exposes both a \emph{shape} (tensor /
set / closure) and a \emph{measurement primitive} used by the cross-
cutting metrics of \S3.

\paragraph{Honest framing.}
The 9-layer count is a \emph{design} decomposition, not an empirical
measurement: every layer has at least one shipped metric
(\texttt{MEASURED} per \texttt{lambda\_layer\_metrics.md}), but the
catalogue itself is the author's taxonomy and is therefore labelled
\DESIGN{}.

% ---------------------------------------------------------------------------
% S2 -- Hyperparameters: full search budget + CFM budget + REINVENT4 + weights
% ---------------------------------------------------------------------------

\subsection*{S2. Hyperparameters: search budget, CFM budget, REINVENT4 config, scoring weights}
\label{supp:s2-hparams}

\paragraph{Full search budget.}
The default $\lambda$-only baseline budget is $N{=}10$ candidates per
pocket, $B{=}60$ rewrite depth, $s{=}1$ and $s{=}5$ MCTS simulation
counts, $k_{\text{every}}{=}1$ prior cadence, and 3 seeds per pocket.
These are the defaults written into
\texttt{molmetal/scripts/r4\_lambda\_only\_run.py} and are identical to
the WF-Lambda-1 protocol \cite{footnote:wf_lambda1}.  The
3-seed$\times$$N{=}10$ pilot results live under
\texttt{molmetal/reports/wf\_lambda1\_N10\_3seed\_v2/} and are summarised
in \S5 of the manuscript.

\paragraph{CFM (Conditional Flow Matching) budget.}
The Round-10 Pt / CFG / Vina end-to-end micro-bench \cite{footnote:round10_e2e}
used hidden\_dim$=32$, $n_{\text{layers}}{=}2$, lr$=2{\times}10^{-4}$,
context\_dropout$=0.1$, $30$ CFM training steps, $8$ ODE integration
steps, $20$ molecules per setting, and Vina exhaustiveness $2$ with
$n_{\text{poses}}{=}1$.  These are the explicit numbers in
\texttt{molmetal/reports/round10\_e2e\_pt\_cfg\_vina.md} (lines
8--17) and are the only CFM configuration under which a \MEASURED{}
Pt-micro-bench currently exists.

\paragraph{REINVENT4 multiproperty configuration.}
The AMD-tuned multiproperty JSON used by the $r_{\text{reinvent4}}$
channel is \texttt{molmetal/configs/reinvent\_multiproperty\_amd.json}
(weights, transforms, and reference molecules; see
\texttt{footnote:wf\_extra2}).  The corresponding end-to-end batch
smoke test lives under
\texttt{molmetal/reports/wf\_extra2\_batch/} (file
\texttt{final.md} summarises the 10-SMILES drug-like batch results and
the offline graceful-fallback behaviour).

\paragraph{Scoring weights.}
The RewardAggregator combines Vina, QuickVina2 parity, MCTS-UCT,
homotype-diversity, pIC50 (calibrated ridge) and REINVENT4 channels
with weights
\begin{equation}
  \label{eq:reward-weights}
  w = (w_{\text{vina}}, w_{\text{qv2}}, w_{\text{mcts}}, w_{\text{homotype}}, w_{\text{pic50}}, w_{\text{reinvent4}})
  = (1.0,\; 0.0,\; 0.3,\; 0.2,\; 0.5,\; 0.0),
\end{equation}
matching the per-channel weights hard-coded in
\texttt{molmetal/molmetal/eval/scoring.py:RewardAggregator} (verified
against the WF-Extra-2 wire report).  Channels with weight $0.0$ are
\emph{off by default} and are only enabled by the corresponding
\texttt{--enable-*} CLI flag.

\paragraph{Honest framing.}
All numbers in this section are \DESIGN{} (defaults pulled from the
canonical scripts/configs); the only \MEASURED{} cell is the CFM
hyperparameter block, which comes from a single-pocket micro-bench and
is therefore \emph{not} a Round-12/13 acceptance number.

% ---------------------------------------------------------------------------
% S3 -- Per-pocket results (placeholder for Round-12 + Round-13 data)
% ---------------------------------------------------------------------------

\subsection*{S3. Per-pocket results}
\label{supp:s3-per-pocket}

\paragraph{Status.}
\DESIGN{}.  This section is a placeholder for the Round-12
$N{=}10{\times}3$ scientific pilot \cite{footnote:round12_pilot} and the
Round-13 100-pocket $\times$ 3-seed full sweep \cite{footnote:round13_sweep}.
Both workstreams are currently pending (WF-3 and WF-5 in the project
tracker; see \texttt{TODO/pending/20\_post\_r10\_r11\_action\_plan.md}).

\paragraph{Planned layout (to be filled in once the experiments ship).}
\begin{itemize}
  \item \textbf{Per-pocket table.}  Columns: pocket-id, pocket-PDB path,
        seed, $N$ candidates requested, $N$ docked, mean Vina score
        (kcal/mol, $\pm$ SD), median Vina, PoseBusters pass-rate, MCTS
        $w_{\text{UCT}}$ at convergence, lambda-only diversity
        (homotype + Tanimoto), pIC50 calibrated-ridge posterior
        (mean $\pm 1\sigma$), REINVENT4 multiproperty score (or
        \texttt{offline-fallback}).
  \item \textbf{Aggregate rollups.}  Pooled mean / 95\% CI over
        pockets, broken down by setting
        (B=60 s=1 vs B=60 s=5; metal-seed on/off).  These are the
        acceptance numbers; per-seed variance is reported separately.
  \item \textbf{Delta vs ablation.}  Signed $\Delta$ between main arm
        and each of the three ablations (CuAAC-only, no-metal-seed,
        no-MCTS), with pooled SE and bootstrap 95\% CI.
\end{itemize}

\paragraph{Honest framing.}
Until the Round-12 and Round-13 numbers ship, this section contains no
\MEASURED{} cell.  Any per-pocket number quoted in the main manuscript
body is drawn directly from the round-10 micro-bench
\cite{footnote:round10_e2e} and is therefore tagged
\MEASURED{}-single-pocket at best.

% ---------------------------------------------------------------------------
% S4 -- Ablation details: CuAAC-only vs all-5
% ---------------------------------------------------------------------------

\subsection*{S4. Ablation details: CuAAC-only vs all-5 click rules}
\label{supp:s4-ablation}

\paragraph{Setup.}
The WF-Lambda-1 ablation harness
(\texttt{molmetal/scripts/r4\_lambda\_only\_run.py}
\texttt{--click-rules cuaac} vs the default 5-rule registry)
generates two arms from the same seed and same $B, s, N$ budget.  The
5-rule registry lives in
\texttt{molmetal/molmetal\_lam/lam\_chem/rules.py} as
$\{r_1, \dots, r_5\} = \{\text{CuAAC, SPAAC, Thiol-ene, Suzuki-Miyaura,
Amide coupling}\}$.

\paragraph{Measured diversity contribution.}
From the WF-Lambda-1 build report
(\texttt{molmetal/reports/wf\_lambda1\_build.md}), the four non-CuAAC
rules together account for approximately 80\% of the per-cell
\emph{structural diversity} contribution as measured by the homotype
distance metric (see \S5 below); CuAAC alone covers the remaining
$\sim$20\%.  This is the empirical justification for retaining all
five rules in the default registry: dropping any of the four
non-CuAAC rules would disproportionately compress the productive space.

\paragraph{Honest framing.}
The 80\% / 20\% split is \MEASURED{} on the WF-Lambda-1 $N{=}10$
3-seed pilot set; it is \emph{not} an extrapolation to the full
100-pocket Round-13 sweep.  The CuAAC-only ablation arm is
\DESIGN{}-in-spirit (no per-pocket acceptance number is reported in
the main text) and is documented here to make the ablation transparent
to reviewers.

% ---------------------------------------------------------------------------
% S5 -- Homotype diversity derivation
% ---------------------------------------------------------------------------

\subsection*{S5. Homotype diversity: derivation and 10-molecule scatter}
\label{supp:s5-homotype}

\paragraph{Derivation.}
The \emph{homotype diversity} metric
$d_{\text{homotype}}(m_1, m_2) \in [0, 1]$ is defined in
\texttt{molmetal/reports/wf\_lambda2\_metric.md} and implemented in
\texttt{molmetal/molmetal\_lam/eval/homotype\_diversity.py}.  It
combines (i) typed-variable overlap (atom-type, hybridisation, ring
membership) under the enriched vocabulary of WF-Lambda-2E
\cite{footnote:wf_lambda2e} and (ii) typed-bond overlap (bond order,
aromatic flag, in-ring flag).  The metric is registered as the 7th
cross-cutting metric in \texttt{r4\_lambda\_only\_run.py}.

\paragraph{Enriched vocabulary.}
The WF-Lambda-2E extension enriches the typed-variable alphabet from
$\{$atom-type, ring-membership$\}$ to
$\{$atom-type, hybridisation, ring-membership, charge-state$\}$.  The
comparison vs the legacy vocabulary is reported in
\texttt{molmetal/reports/wf\_lambda2e\_extend.md}, with the enriched
vocabulary yielding a tighter correlation between homotype distance
and downstream oracle reward on the 10-mol test set
(\texttt{wf\_lambda2e\_compare/final.md}).

\paragraph{10-molecule scatter plot.}
The 10-molecule scatter plot referenced in \S5 of the main text is
produced by \texttt{molmetal/scripts/wf\_lambda2e\_scatter.py} (input:
\texttt{molmetal/reports/wf\_lambda2e\_compare/test\_set.smi};
output: \texttt{wf\_lambda2e\_compare/scatter.png}).  Each point is one
generated molecule; axes are homotype distance ($x$) vs Tanimoto
distance to the seed ($y$), coloured by rule-of-origin.

\paragraph{Honest framing.}
The 10-mol test set is a \emph{convenience panel}, not a hold-out; the
metric itself is \DESIGN{} (definition) and the 10-mol scatter is
\MEASURED{} (computed on the convenience panel).  The claim that
homotype distance is a useful diversity proxy requires the full
Round-13 sweep and is therefore tagged \DESIGN{} in the main text.

% ---------------------------------------------------------------------------
% S6 -- pIC50 conditioned-ridge baseline
% ---------------------------------------------------------------------------

\subsection*{S6. pIC50 conditioned-ridge baseline}
\label{supp:s6-pic50}

\paragraph{Configuration.}
The pIC50 baseline is a per-assay ridge regression with a SMILES
Morgan-fingerprint input ($\text{radius}=2$, $\text{nBits}=2048$).  The
training pipeline is the WF-Extra-1 workstream
\cite{footnote:wf_extra1_full} (3-seed $\times$ 50-epoch retrain with
a censor-aware loader).  Calibration is performed per-assay: each
ChEMBL target family gets its own ridge weight $\beta_t$ and bias
$\alpha_t$, fit on the censored training pairs
$(\text{Mol}, \text{pIC50})$ with bound-aware accuracy as the held-out
metric.

\paragraph{Bound-aware accuracy.}
Censored measurements (``$>$X'' or ``$<$X'') are not dropped; they
contribute a one-sided hinge loss.  The exact loss form and the
censored-loader implementation are in
\texttt{molmetal/reports/wf\_extra1\_censor\_loader.md} and
\texttt{molmetal/molmetal/eval/calibrate\_conditioned\_pic50\_baseline.py}.

\paragraph{Inputs to the RewardAggregator.}
The $r_{\text{pic50}}$ channel returns the calibrated ridge posterior
mean, clipped to $[0, 10]$ (pIC50 is by convention in $[0, 10]$ on
ChEMBL).  The default weight is $w_{\text{pic50}} = 0.5$ (see
Eq.~\ref{eq:reward-weights}).  When the ridge is unavailable (e.g.\ no
assay matches), the channel gracefully returns $0.0$.

\paragraph{Honest framing.}
The ridge hyperparameters are \DESIGN{} (cross-validated defaults);
the bound-aware accuracy and the 3-seed $\times$ 50-epoch retrain
numbers are \MEASURED{} per \texttt{wf\_extra1\_full/final.md}.

% ---------------------------------------------------------------------------
% S7 -- REINVENT4 multiproperty scoring protocol
% ---------------------------------------------------------------------------

\subsection*{S7. REINVENT4 multiproperty scoring protocol}
\label{supp:s7-reinvent4}

\paragraph{Configuration file.}
The AMD-tuned multiproperty JSON
(\texttt{molmetal/configs/reinvent\_multiproperty\_amd.json}) is the
single source of truth for the REINVENT4 scoring channel.  It encodes:
\begin{itemize}
  \item Per-component weights $w_i$ (QED, SA, logP, pIC50-proxy, BBB).
  \item Per-component transforms (sigmoid, clip, log).
  \item Reference molecules for the SA / logP calibration.
\end{itemize}

\paragraph{Wiring.}
The subprocess adapter lives in
\texttt{molmetal/molmetal/adapters/reinvent4\_subprocess\_adapter.py}
and is invoked by the \texttt{r\_reinvent4} channel of the
RewardAggregator.  The end-to-end wire report is
\texttt{molmetal/reports/wf\_extra2\_wire.md}; the 10-SMILES
drug-like batch smoke test is
\texttt{molmetal/reports/wf\_extra2\_batch/final.md}.

\paragraph{Offline graceful fallback.}
When the REINVENT4 backend is not reachable (e.g.\ no licence, no
network), the channel returns $0.0$ and logs a one-line warning;
the rest of the RewardAggregator keeps running.  This behaviour is
unit-tested in \texttt{tests/test\_reinvent4\_channel.py} (see
\texttt{wf\_extra2\_wire.md} for the test catalogue).

\paragraph{Honest framing.}
The configuration is \DESIGN{} (manually tuned on a ChEMBL drug-like
sample).  The 10-SMILES batch test is \MEASURED{}.  The full
REINVENT4 scoring run on the Round-13 sweep is \DESIGN{} (to be
filled in once WF-5 ships).

% ---------------------------------------------------------------------------
% S8 -- Vina vs QuickVina2 parity protocol + scatter
% ---------------------------------------------------------------------------

\subsection*{S8. Vina 1.2.7 vs QuickVina 2 engine-parity protocol}
\label{supp:s8-vina-qv2}

\paragraph{Setup.}
The Round-11 parity study
(\texttt{molmetal/reports/round11\_engine\_parity\_n50.md},
\cite{footnote:round11_parity}) docked the same 50 cross-docked
receptor-ligand pairs through both engines (Vina 1.2.7 and QuickVina 2)
under matched exhaustiveness and seed settings.  Per-pocket parity is
quantified by Pearson $r$ and by the 95\% limits of agreement on the
docked-score axis.

\paragraph{Result.}
The reported correlation is $r = 0.9983$ over $N{=}50$ pairs, with
bias $< 0.1$ kcal/mol and 95\% limits of agreement $\pm 0.6$
kcal/mol.  The per-pocket scatter plot is
\texttt{molmetal/reports/round11\_engine\_parity/scatter.png}.

\paragraph{Interpretation.}
The two engines are statistically interchangeable at the precision of
this paper; the manuscript therefore reports a single Vina-or-QuickVina2
score per pocket and tags the source engine alongside the number.

\paragraph{Honest framing.}
The $r = 0.9983$ is \MEASURED{} on $N{=}50$ pairs; the conclusion that
the two engines are interchangeable for the purposes of this work is
\DESIGN{} (based on the magnitude of the bias and limits of agreement
relative to the manuscript's success-criterion threshold of
$|\Delta| \geq 0.2$ kcal/mol).

% ---------------------------------------------------------------------------
% S9 -- beta-NF strong normalisation proof
% ---------------------------------------------------------------------------

\subsection*{S9. $\beta$-normal-form strong normalisation}
\label{supp:s9-betanf}

\paragraph{Statement.}
For the typed fragment of the Molecular Lambda Calculus restricted to
atoms, bonds, applications and the three combinators of
\texttt{molmetal/molmetal\_lam/lam\_chem/rules.py}, every
well-typed term has a (unique) $\beta$-normal form, and the
normalisation is \emph{strongly normalising} (every reduction
sequence terminates in the same normal form).  The proof is given in
Appendix~B (\texttt{paper/appendices/betanf\_semantics.tex}, see also
\texttt{WF-Lambda-4}).

\paragraph{Sketch.}
The proof proceeds by induction on the size of the typed term, using
the fact that (i) every combinator reduces in a constant number of
steps, (ii) every click-rule product is closed under the
valence-BNF predicate of \texttt{molmetal/molmetal\_lam/lam\_chem/bnf.py},
and (iii) the rewrite system is locally confluent and terminating at
each step.  The full proof is the content of
\texttt{paper/appendices/betanf\_semantics.tex}.

\paragraph{Honest framing.}
The statement is \DESIGN{}; the proof is \DESIGN{} (typed fragment,
not the full untyped $\lambda$-calculus).  No empirical claim is made
in this section.

% ---------------------------------------------------------------------------
% S10 -- Closure theorem + 5-click productive space algorithm
% ---------------------------------------------------------------------------

\subsection*{S10. Closure theorem and 5-click productive-space algorithm}
\label{supp:s10-closure}

\paragraph{Theorem.}
See Appendix~A (\texttt{paper/appendices/closure\_theorem.tex},
\texttt{WF-Lambda-4 audit + closure-theorem surface},
\texttt{molmetal/reports/wf\_lambda4\_closure\_theorem.md}).
The closure theorem states that the BFS-closure
$\mathcal{P}(\Sigma, \mathcal{R}, B)$ of any seed term under the five
click rules is a well-formed, RDKit-sanitisable, valence-BNF-satisfied
\emph{constructive synthesis space}, and that every term in
$\mathcal{P}$ admits a $\beta$-reduction witness path back to the seed.

\paragraph{5-click productive-space algorithm.}
\begin{enumerate}
  \item \textbf{Seed.}  Choose an initial closed term $S \in \Sigma$
        (e.g.\ cisplatin, $[\text{Pt}(\text{NH}_3)_2\text{Cl}_2]$).
  \item \textbf{Initialise.}  Set frontier $F \leftarrow \{S\}$,
        visited $V \leftarrow \emptyset$, depth $d \leftarrow 0$.
  \item \textbf{Expand.}  For each term $t \in F$ and each rule
        $r_i \in \mathcal{R}$, enumerate matching educts
        $(a, b)$ with $a \in V \cup F$, $b \in \Sigma$, and apply
        $r_i$ to obtain the product $c = r_i.\text{reduce}(a, b)$.
  \item \textbf{Validate.}  Drop any $c$ that fails the valence-BNF
        predicate or the RDKit sanitiser.
  \item \textbf{Update.}  Set $V \leftarrow V \cup F$, $F \leftarrow
        \{c \text{ not in } V\}$, $d \leftarrow d+1$.
  \item \textbf{Stop.}  Halt when $d = B$ or $F = \emptyset$.
\end{enumerate}
The algorithm is implemented in
\texttt{molmetal/molmetal\_lam/closure.py} and is the back-end of the
\texttt{r\_mcts} channel of the RewardAggregator.

\paragraph{Honest framing.}
The theorem statement and the algorithm are \DESIGN{}.  The
implementation is \DESIGN{}-with-\texttt{MEASURED}-unit-tests: the
BFS termination, BNF validity, and RDKit sanitisation are covered by
\texttt{tests/test\_closure.py}, but no large-scale empirical study of
the productive-space size has been published at the time of writing.

% ---------------------------------------------------------------------------

\paragraph{Reproducibility.}
Every supplementary item above points to a file under
\texttt{molmetal/reports/}, \texttt{molmetal/configs/},
\texttt{molmetal/molmetal\_lam/}, or \texttt{paper/appendices/}.  The
reproduction environment (\texttt{uv}-managed Python 3.12, ROCm 7.2,
triton-rocm 3.8.0, RX 7800 XT gfx1101 wave64) is documented in
\texttt{molmetal/reports/wf\_extra2\_wire.md} appendix B and is the
same environment as the main text.
